\documentclass[12pt]{article}
\usepackage[utf8]{inputenc}
\usepackage[T1]{fontenc}
\usepackage{hyperref}
\title{Sagalla Language Resources for K-12 Curriculum Development}
\author{Compiled by Claude Sonnet 5 (claude-sonnet-5) for Andr\'as Kornai}
\date{}
\begin{document}
\maketitle

\section{Introduction}
Sagalla (Kisagalla; ISO 639-3 tga; Glottocode saga1262) is a Bantu language of the Taita Hills, Taita-Taveta County, Kenya. Sources disagree on its classification and population: one treats Saghala as a dialect of Taita, and the speaker figures range from 8,000 to 100,000 or more. Published material is sparse.

\section{Primary Educational Resources}
No primers, textbooks or school programmes were found in the search results.

\section{Linguistic and Reference Materials}
\begin{itemize}
\item Wikipedia overview of the Sagalla language: \url{https://en.wikipedia.org/wiki/Sagalla_language}
\item University of Florida Africa document describing printed Taita materials (a Bible, hymn book, folklore pamphlet, proverbs and devotional) and a pending bilingual Taita-English dictionary project: \url{https://africa.ufl.edu/wp-content/uploads/sites/167/page-21_lugano.pdf}
\item Endangered Languages Project entry: \url{https://endangeredlanguages.com/node/96952}
\item Omniglot page on Taita (treats Saghala as a dialect): \url{https://www.omniglot.com/writing/taita.htm}
\end{itemize}

\section{Religious and Community Resources}
\begin{itemize}
\item Scripture Earth page for tga (New Testament): \url{https://www.scriptureearth.org/00eng.php?iso=tga}
\item Joshua Project language page: \url{https://www.joshuaproject.net/index.php/languages/tga}
\item Joshua Project people-group page, Sagalla Taita in Kenya: \url{https://www.joshuaproject.net/index.php/people_groups/14635}
\end{itemize}

\section{Digital and Audio Resources}
The New Testament is reported as available through Bible.is and YouVersion (via Scripture Earth and Joshua Project). The search results did not show whether it includes audio, and no other digital resources were found.

\section{Curriculum Development Context}
The language is not known to be taught in schools. No dictionary was found; the Taita-English dictionary project was reported as pending. A curriculum would have to build on the New Testament text and community work, and the Sagalla/Taita relationship would need clarifying first.
\end{document}